\def\handoutsay{false}\def\handoutfill{false}
\documentclass[12pt]{article}
\newcommand{\handoutrunning}{Lecture 12: Linear Models II}
\providecommand{\handoutsay}{true}
\usepackage[T1]{fontenc}
\usepackage{lmodern,microtype}
\usepackage[a4paper,margin=15mm,headheight=14pt,headsep=5mm]{geometry}
\usepackage{amsmath,amssymb,amsthm,mathtools}
\usepackage{enumitem,booktabs,needspace}
\usepackage{pifont}
\usepackage[dvipsnames]{xcolor}
\usepackage{fancyhdr}
\usepackage{ifthen}
\usepackage[hidelinks]{hyperref}
\hypersetup{pdfauthor={Chang-Ye Tu}}
\pagestyle{fancy}
\fancyhf{}
\fancyhead[L]{\small Linear Algebra}
\fancyhead[R]{\small\handoutrunning}
\fancyfoot[C]{\small\thepage}
\renewcommand{\headrulewidth}{0.3pt}
\setlength{\parindent}{0pt}
\setlength{\parskip}{4pt plus 1pt minus 1pt}
\setlength{\emergencystretch}{1.5em}
\usepackage{etoolbox}
\newcommand{\handoutdisplayskips}{\setlength{\abovedisplayskip}{7pt plus 2pt minus 2pt}%
  \setlength{\belowdisplayskip}{7pt plus 2pt minus 2pt}}
\apptocmd{\normalsize}{\handoutdisplayskips}{}{\PackageWarning{handout}{could not patch normalsize}}
\handoutdisplayskips
\setlist[enumerate]{label=(\alph*),leftmargin=*,itemsep=2pt,topsep=3pt}
\setlist[itemize]{leftmargin=*,itemsep=2pt,topsep=3pt}
\newcommand{\R}{\mathbb R}
\newcommand{\C}{\mathbb C}
\newcommand{\iu}{\mathrm{i}}
\newcommand{\contradiction}{\mathord{\text{\large\ding{56}}}}
\newcommand{\one}{\mathbf 1}
\DeclareMathOperator{\tr}{tr}
\DeclareMathOperator{\col}{col}
\DeclareMathOperator{\diag}{diag}
\DeclareMathOperator{\rank}{rank}
\DeclareMathOperator\prb{\mathsf{P}}
\DeclareMathOperator\expc{\mathsf{E}}
\DeclareMathOperator\var{var}
\DeclareMathOperator\cov{cov}
\DeclareMathOperator\cor{cor}
\DeclareMathOperator\argmin{argmin}
\DeclareMathOperator\argmax{argmax}
\DeclareMathOperator{\Span}{span}
\theoremstyle{definition}
\newtheorem{theorem}{Theorem}
\newtheorem{proposition}[theorem]{Proposition}
\newtheorem{lemma}[theorem]{Lemma}
\newtheorem{corollary}[theorem]{Corollary}
\newtheorem{definition}[theorem]{Definition}
\newtheorem{example}[theorem]{Example}
\newtheorem{remark}[theorem]{Remark}
\newtheorem{fact}{Fact}
\renewcommand{\thefact}{F\arabic{fact}}
\AddToHook{cmd/subsection/before}{\par\Needspace{9\baselineskip}}
\AddToHook{env/theorem/before}{\par\Needspace{4\baselineskip}}
\AddToHook{env/proposition/before}{\par\Needspace{4\baselineskip}}
\AddToHook{env/lemma/before}{\par\Needspace{4\baselineskip}}
\AddToHook{env/definition/before}{\par\Needspace{4\baselineskip}}
\AddToHook{env/example/before}{\par\Needspace{4\baselineskip}}
\definecolor{saybg}{gray}{0.86}
\NewDocumentEnvironment{say}{+b}
  {\ifthenelse{\equal{\handoutsay}{true}}{\par\Needspace{2\baselineskip}\noindent\colorbox{saybg}{\parbox{\dimexpr\linewidth-2\fboxsep\relax}%
   {\small\sffamily\ding{43}\enspace\textbf{Say.}\enspace#1}}\par}{}}
  {}
\providecommand{\handoutfill}{false}
\newsavebox{\boardbox}
\newlength{\boardht}
\definecolor{boardrule}{gray}{0.55}
\makeatletter
\newcommand{\boardstart}{\par\if@inlabel\leavevmode\par\fi}
\makeatother
\NewDocumentEnvironment{board}{o +b}
  {\boardstart\IfValueT{#1}{\noindent#1\par\nopagebreak}%
   \noindent\begin{lrbox}{\boardbox}\begin{minipage}{\linewidth}#2\end{minipage}\end{lrbox}%
   \setlength{\boardht}{\dimexpr\ht\boardbox+\dp\boardbox\relax}%
   \ifthenelse{\equal{\handoutfill}{true}}%
     {\setlength{\boardht}{1.4\boardht}%
      {\setlength{\fboxsep}{0pt}\setlength{\fboxrule}{0.4pt}\color{boardrule}%
       \framebox[\dimexpr\linewidth-2\fboxrule\relax]{\rule{0pt}{\boardht}}}}%
     {\ifthenelse{\equal{\handoutsay}{true}}%
        {\llap{\color{boardrule}\rule[-\dp\boardbox]{1.5pt}{\boardht}\hspace{4pt}}}{}%
      \usebox{\boardbox}}%
   \par}
  {}
\newcommand{\blank}[1]{\ifthenelse{\equal{\handoutfill}{true}}{\underline{\phantom{#1}}}{#1}}
\newcommand{\lec}[1]{Lecture~#1}
\newcommand{\handouttitle}[2]{%
  \begin{center}
  {\LARGE\sffamily\bfseries #2\par}\vspace{3pt}
  {\sffamily Linear Algebra \textperiodcentered\ \lec{#1}\par}
  \ifthenelse{\equal{\handoutsay}{true}}{\vspace{3pt}{\small\sffamily\colorbox{saybg}{\ding{43}\enspace\textbf{Say.}}\enspace
    teacher's copy: the shaded boxes are not in the students' copy.\par}}{}
  \ifthenelse{\equal{\handoutfill}{true}}{\vspace{3pt}{\small\sffamily Fill-in copy: complete the framed spaces in
    class.\par}}{}
  \end{center}}

\hypersetup{pdftitle={Linear Algebra, Lecture 12: Linear Models II}}
\newcommand{\boardtop}{\vspace*{-\abovedisplayskip}}
\DeclareMathOperator{\SE}{SE}
\begin{document}
\handouttitle{12}{Linear Models II}

\textbf{Last time} (\lec{11}): \(\beta^0=GX'y\) (\(G\) a g-inverse of \(X'X\)) solves the normal equations;
\(H=XGX'\) does not depend on \(G\); \(\hat\sigma^2=y'My/(n-r)\) (\(M=I_n-H\)) is unbiased. Earlier results are cited where they are
used; F1--F3: probability, no proof. \(e\) is the vector of residuals; unit vectors are columns of \(I\).

\begin{enumerate}[label=\textbf{F\arabic*.},ref=F\arabic*,leftmargin=*,itemsep=0pt,parsep=0pt,topsep=2pt]
\item\label{f-prob} \(\prb(u\notin\mathcal A)=1-\prb(u\in\mathcal A)\); \(u_1,\dots,u_m\) are \textbf{independent} if
\(\prb(u_k\in\mathcal A_k\ \forall k)=\prod_k\prb(u_k\in\mathcal A_k)\); then so are \(g_1(u_1),\dots,g_m(u_m)\).
\item\label{f-chisq} \(\chi^2_r\), \(t_\nu\), \(F_{a,b}\) (\textbf{degrees of freedom} \(r,\nu,a,b\geqslant1\)): the
distributions of \(\sum_{i\leqslant r}w_i^2\), \(Z/\sqrt{Y/\nu}\), \((Y_1/a)/(Y_2/b)\) for independent \(w_i,Z\sim N(0,1)\),
\(Y\sim\chi^2_\nu\), \(Y_1\sim\chi^2_a\), \(Y_2\sim\chi^2_b\).
\item\label{f-quantile} For \(u\) as in \ref{f-chisq}: \(\prb(u=c)=0\); the \textbf{quantile} \(t_{\nu,\tau}\),
\(F_{a,b,\tau}\) (\(0<\tau<1\)) is the unique \(c\) with \(\prb(u\leqslant c)=\tau\).
\end{enumerate}

\subsection*{1. Estimable functions and their best linear unbiased estimators}

\begin{definition}\label{h-def-estimable}
In the linear model \(y=X\beta+\varepsilon\) (\lec{11}, Definition~1: \(X\in\R^{n\times p}\) known, \(\beta\in\R^p\)
unknown, \(\varepsilon\in\R^n\) random), with \(X\) of any rank, assume (A1) \(\expc(y)=X\beta\), and let \(a\in\R^p\).
A \textbf{linear estimator} is a random variable \(c'y\) with a constant \(c\in\R^n\); it is a \textbf{linear unbiased
estimator} of \(a'\beta\) if \(\expc(c'y)=a'\beta\) for every \(\beta\in\R^p\). \(a'\beta\) is \textbf{estimable} if it
has a linear unbiased estimator. A linear unbiased estimator of \(a'\beta\) with the smallest variance among all of them
is a \textbf{best linear unbiased estimator} (\textbf{BLUE}) of \(a'\beta\).
\end{definition}

\Needspace{26\baselineskip}
\begin{theorem}\label{h-thm-estimable}
Assume (A1), and let \(a\in\R^p\).
\begin{enumerate}
\item A linear estimator \(c'y\) is unbiased for \(a'\beta\) \(\iff\) \(X'c=a\). Hence \(a'\beta\) is estimable \(\iff\)
\(a\in\mathcal C(X')\) \(\iff\) \(a\perp\mathcal N(X)\). If \(\rank X=p\), every \(a'\beta\) is estimable.
\item Let \(G\) be a g-inverse of \(X'X\) (\lec{11}, Definition~8: \(AGA=A\) for \(A=X'X\)) and \(H=XGX'\). If
\(a=X'c\), then \(a'\hat\beta=c'\hat y=c'Hy\) (\(\hat y=X\hat\beta\)) for every least-squares estimate \(\hat\beta\), a
minimizer of \(\lVert y-X\beta\rVert^2\) (\lec{11}, Definition~1): the same linear estimator for all of them; and
\(\expc(a'\hat\beta)=a'\beta\).
\item If \(a'\beta\) is not estimable, then \(a'\hat\beta\) takes every real value as \(\hat\beta\) runs through the
least-squares estimates.
\end{enumerate}
\end{theorem}
\begin{proof}
\begin{board}[(a) \(\expc(c'y)=c'\expc(y)=c'X\beta\) (\lec{9}, Theorem~3(b): \(\expc(Az+b)=A\expc(z)+b\); (A1));
\(z\perp W\) means \(z'w=0\) for all \(w\in W\) (\lec{6}, Definition~4); \(\rank A=\dim\mathcal C(A)\) (\lec{4},
Definition~10); subspaces \(U\subseteq W\) with \(\dim U=\dim W\) are equal (\lec{4}, Theorem~5(a)):]
\boardtop
\begin{align*}
c'X\beta&=a'\beta\ (\text{all }\beta)\implies c'X\col_j(I_p)=a'\col_j(I_p)\ (\text{all }j)&&\text{(\lec{1}, Proposition~11(a))}\\
&\implies c'X=a'\iff X'c=a&&\text{(\lec{1}, Theorem~8)}\\
X'c&=a\implies\expc(c'y)=c'X\beta=a'\beta\\
&\implies\bigl(a'\beta\text{ estimable}\iff a\in\{X'c\}=\mathcal C(X')\bigr)&&\text{(\lec{4}, Definition~7)}\\
\mathcal N(X)&=\mathcal C(X')^\perp\implies\mathcal C(X')\subseteq\mathcal N(X)^\perp&&\text{(\lec{6}, Corollary~6)}\\
\dim\mathcal C(X')&=\rank X'=\rank X=p-\dim\mathcal N(X)&&\text{(\lec{4}, Theorems~12(a), 11)}\\
&=\dim\mathcal N(X)^\perp\implies\mathcal C(X')=\mathcal N(X)^\perp&&\text{(\lec{6}, Theorem~5(c))}\\
\rank X&=p\implies\rank X'=p\implies\mathcal C(X')=\R^p&&\text{(\lec{4}, Theorem~12(a),(c))}
\end{align*}
\end{board}
\begin{board}[(b) \(a=X'c\); \(\expc(c'Hy)=c'H\expc(y)\) (\lec{9}, Theorem~3(b)) and (A1); all least-squares estimates
give the same \(\hat y=X\hat\beta\) (\lec{11}, Theorem~3(b)):]
\boardtop
\begin{align*}
X\hat\beta&=\hat y=Hy\implies a'\hat\beta=c'X\hat\beta=c'\hat y=c'Hy&&\text{(\lec{11}, Theorem~10(e))}\\
\expc(c'Hy)&=c'HX\beta=c'X\beta=a'\beta&&\text{(\lec{11}, Theorem~10(c))}
\end{align*}
\end{board}
\begin{board}[(c) By (a), \(a\notin\mathcal N(X)^\perp\); \(X'X\beta=X'y\) has a solution \(\hat\beta\) (\lec{11},
Theorem~2(a)); the least-squares estimates are the solutions of \(X'X\beta=X'y\) (\lec{11}, Theorem~3(a)):]
\boardtop
\begin{align*}
a'z&\ne0\ \text{for some }z\in\mathcal N(X)\\
Xz&=0\implies X'X(\hat\beta+tz)=X'X\hat\beta=X'y\quad(t\in\R)\implies\hat\beta+tz\text{ is a least-squares estimate}\\
a'(\hat\beta+tz)&=a'\hat\beta+t\,a'z\ \text{takes every real value}
\end{align*}
\end{board}
\nopagebreak
\end{proof}

\Needspace{14\baselineskip}
\begin{example}\label{h-exm-estimable}
In \lec{11}, Example~7, the one-way model \(y_i=\mu+\alpha_l+\varepsilon_i\) of \lec{11}, Theorem~6 has
\(\beta=(\mu,\alpha_1,\alpha_2,\alpha_3)'\), \(y=(5,7,6,3,5,2)'\) in groups of sizes \((n_1,n_2,n_3)=(3,2,1)\),
\(X'X=\left(\begin{smallmatrix}6&3&2&1\\3&3&0&0\\2&0&2&0\\1&0&0&1\end{smallmatrix}\right)\), \(X'y=(28,18,8,2)'\) and the
least-squares estimates \((-t,6+t,4+t,2+t)'\) (\(t\in\R\)). Which of \(\mu+\alpha_1\), \(\alpha_1-\alpha_2\), \(\mu\),
\(\alpha_1\) are estimable? With the group means \(\bar y_1=\frac13(y_1+y_2+y_3)\), \(\bar y_2=\frac12(y_4+y_5)\),
\(\bar y_3=y_6\), find the estimator of \(\alpha_1-\alpha_2\) of Theorem~\ref{h-thm-estimable}(b).
\end{example}
\textit{Solution.}
\begin{board}[\(\mathcal N(X)=\Span\{(-1,1,1,1)'\}\) (\lec{11}, Theorem~6(b)); Theorem~\ref{h-thm-estimable}(a),(c):]
\boardtop
\begin{align*}
&a_0\mu+a_1\alpha_1+a_2\alpha_2+a_3\alpha_3\ \text{is estimable}\iff-a_0+a_1+a_2+a_3=0\\
&\mu+\alpha_1:\ -1+1=0,\quad\alpha_1-\alpha_2:\ 1-1=0\implies\text{estimable; on }(-t,6+t,4+t,2+t)'\text{: values }6,\ 2\\
&\mu:\ -1\ne0,\quad\alpha_1:\ 1\ne0\implies\text{not estimable; values }{-t},\ 6+t\text{ change with }t
\end{align*}
\end{board}
\begin{board}[\(X'y\) holds the total and the group totals \(\col_j(X)'y\) (\lec{1}, Theorem~4(d), Definition~7);
\lec{11}, Theorem~3(a):]
\boardtop
\begin{align*}
X'X(0,\bar y_1,\bar y_2,\bar y_3)'&=(3\bar y_1+2\bar y_2+\bar y_3,\ 3\bar y_1,\ 2\bar y_2,\ \bar y_3)'=X'y\\
\hat\beta&=(0,\bar y_1,\bar y_2,\bar y_3)'\ \text{is a least-squares estimate}\\
\hat\alpha_1-\hat\alpha_2&=\bar y_1-\bar y_2=c'y,\qquad c=\bigl(\tfrac13,\tfrac13,\tfrac13,-\tfrac12,-\tfrac12,0\bigr)'\\
X'c&=(1-1,\ 1,\ -1,\ 0)'=(0,1,-1,0)'
\end{align*}
\end{board}

\Needspace{17\baselineskip}
\begin{theorem}\label{h-thm-blue}
Assume (A1)--(A3) (\lec{11}, Definition~1: (A2) \(\expc(\varepsilon)=0\), (A3) \(\var(\varepsilon)=\sigma^2I_n\),
\(\sigma^2>0\)), with \(X\) of any rank; let \(G\) be a g-inverse of \(X'X\), \(H=XGX'\), \(M=I_n-H\), and let
\(a'\beta\) be estimable. Then \(a'\hat\beta\), the same for every least-squares estimate \(\hat\beta\), is the BLUE of
\(a'\beta\): it is a linear unbiased estimator of \(a'\beta\), and every linear unbiased estimator \(c'y\) of \(a'\beta\)
has \(\var(c'y)\geqslant\var(a'\hat\beta)\), with equality only if \(c'y=a'\hat\beta\) for every \(y\). Moreover
\(\var(a'\hat\beta)=\sigma^2a'Ga\) (for every such \(G\)).
\end{theorem}
\begin{proof}
\begin{board}[\(c'y\) a linear unbiased estimator of \(a'\beta\); (A3) gives \(\var(y)=\sigma^2I_n\) (\lec{11},
Definition~1); \(\var(Az+b)=A\var(z)A'\) (\lec{9}, Theorem~3(b)); \(H'=H=H^2\), \(M'=M=M^2\) (\lec{11},
Theorem~10(c)); \(\lVert u\rVert^2=u'u\) is \(0\) only for \(u=0\) (\lec{6}, Definition~1); \(u=Mc\):]
\boardtop
\begin{align*}
X'c&=a,\quad a'\hat\beta=c'Hy\ \text{is linear unbiased}&&\text{(Theorem~\ref{h-thm-estimable}(a),(b))}\\
\var(c'y)&=\sigma^2c'c,\quad\var(c'Hy)=\sigma^2c'HH'c=\sigma^2c'Hc=\sigma^2a'Ga\\
\var(c'y)-\var(a'\hat\beta)&=\sigma^2c'(I_n-H)c=\sigma^2c'M'Mc=\sigma^2\lVert u\rVert^2\geqslant0\\
\sigma^2\lVert u\rVert^2&=0\implies u=0\implies c=Hc\implies c'y=c'Hy=a'\hat\beta\ (\text{all }y)
\end{align*}
\end{board}
\nopagebreak
\end{proof}
So \(\var(c'y)=\var(a'\hat\beta)+\sigma^2\lVert Mc\rVert^2\). Under (A4) \(\rank X=p\), every \(a'\beta\) is
estimable (Theorem~\ref{h-thm-estimable}(a)): Theorem~\ref{h-thm-blue} gives \(\var(a'\tilde\beta)\geqslant\var(a'\hat\beta)\)
for linear unbiased \(\tilde\beta\) of \(\beta\) (\lec{11}, Theorem~13).

\Needspace{15\baselineskip}
\begin{example}\label{h-exm-solution}
In Example~\ref{h-exm-estimable}, with the g-inverse \(G=\diag(0,\frac13,\frac12,1)\) of \(X'X\) (\lec{11},
Example~11), under (A1), (A3): find \(\beta^0=GX'y\), \(\expc(\beta^0)\), \(\var(\beta^0)\) and
\(\var(\bar y_1-\bar y_2)\).
\end{example}
\textit{Solution.}
\begin{board}[\(GX'X\): the rows of \(X'X\) times \(0,\frac13,\frac12,1\) (\lec{1}, Proposition~13(a)); \lec{9},
Theorem~3(b):]
\boardtop
\begin{align*}
GX'X&=\begin{pmatrix}0&0&0&0\\1&1&0&0\\1&0&1&0\\1&0&0&1\end{pmatrix},\qquad
\beta^0=GX'y=(0,\bar y_1,\bar y_2,\bar y_3)'\\
\expc(\beta^0)&=GX'X\beta=(0,\ \mu+\alpha_1,\ \mu+\alpha_2,\ \mu+\alpha_3)'&&\text{((A1))}\\
\var(\beta^0)&=GX'(\sigma^2I_n)XG'=\sigma^2GX'XG=\sigma^2\diag(0,\tfrac13,\tfrac12,1)&&\text{((A3); \(G'=G\))}\\
\var(\bar y_1-\bar y_2)&=a'\var(\beta^0)a=\sigma^2\bigl(\tfrac13+\tfrac12\bigr)=\tfrac56\sigma^2=\sigma^2a'Ga&&(a=(0,1,-1,0)')
\end{align*}
\end{board}
So \(\beta^0\) is the \(\hat\beta\) of Example~\ref{h-exm-estimable}, and \(\bar y_l\) estimates \(\mu+\alpha_l\), not
\(\alpha_l\); \(GX'XG\) has the columns of \(GX'X\) times \(0,\frac13,\frac12,1\) (\lec{1}, Proposition~13(a)); and
\(\frac56\sigma^2\) is the variance of Theorem~\ref{h-thm-blue}.

\subsection*{2. The normal linear model; \(t\) intervals and tests}

\begin{definition}\label{h-def-normal}
The \textbf{normal linear model} is the linear model of \lec{11}, Definition~1, with \(X\in\R^{n\times p}\) of any rank,
under (A5) \(\varepsilon\sim N_n(0,\sigma^2I_n)\) for a number \(\sigma^2>0\) (\lec{10}, Definition~1: multivariate
normal). Then (A2), (A3) hold (\lec{10}, Proposition~2:
\(z\sim N_n(\mu,\Sigma)\Rightarrow\expc(z)=\mu\), \(\var(z)=\Sigma\)), and \(y=X\beta+I_n\varepsilon\sim
N_n(X\beta,\sigma^2I_n)\) (\lec{10}, Theorem~3: \(Az+b\sim N_m(A\mu+b,A\Sigma A')\)), so (A1) holds; and
\(y_1,\dots,y_n\) are independent (\ref{f-prob}), \(y_i\sim N(x_i'\beta,\sigma^2)\), \(x_i'\) the row \(i\) of \(X\)
(\lec{10}, Theorem~6(b): \(z\sim N_n(\mu,\diag(\sigma_1^2,\dots,\sigma_n^2))\iff\) the \(z_i\) are independent
\(N(\mu_i,\sigma_i^2)\)). Let \(\alpha\in(0,1)\). A \textbf{hypothesis} \(H_0\) is a condition on the parameters, such as
\(K\beta=m\). A \textbf{test} of \(H_0\) is given by a \textbf{rejection region} \(\mathcal R\), a set of possible values
of \(y\): \(H_0\) is rejected if the observed \(y\) lies in \(\mathcal R\). The test has \textbf{level} \(\alpha\) if
\(\prb(y\in\mathcal R)=\alpha\) for all parameter values that satisfy \(H_0\). Usually \(\mathcal R\) is described by a
\textbf{test statistic} \(T\) computed from \(y\), as \(\{T>c\}\) or \(\{\lvert T\rvert>c\}\). When \(H_0\) is rejected at
level \(\alpha\), the departure from \(H_0\) is \textbf{significant} at level \(\alpha\). A \textbf{confidence interval}
of level \(1-\alpha\) for \(\theta\in\R\) is an interval \([L,U]\) computed from \(y\) with
\(\prb(L\leqslant\theta\leqslant U)=1-\alpha\) for all parameter values (\lec{10}, Definition~18). The distributions
\(\chi^2_r\), \(t_\nu\), \(F_{a,b}\) and their quantiles \(t_{\nu,\tau}\), \(F_{a,b,\tau}\) are those of \ref{f-chisq},
\ref{f-quantile}.
\end{definition}

\Needspace{16\baselineskip}
\begin{theorem}\label{h-thm-sampling}
In the normal linear model let \(\rank X=r<n\), \(G\) a g-inverse of \(X'X\), \(\beta^0=GX'y\), \(H=XGX'\),
\(M=I_n-H\), and let \(K\in\R^{q\times p}\) be such that every entry of \(K\beta\) is estimable.
\begin{enumerate}
\item \(K\beta^0\sim N_q(K\beta,\sigma^2KGK')\), and \(K\beta^0\), \(KGK'\) are the same for every g-inverse \(G\) of
\(X'X\). Under (A4), \(\hat\beta=(X'X)^{-1}X'y\sim N_p(\beta,\sigma^2(X'X)^{-1})\).
\item \(\mathrm{SSE}/\sigma^2\sim\chi^2_{n-r}\), and \(\prb(\mathrm{SSE}=0)=0\).
\item \(K\beta^0\) and \(\mathrm{SSE}\) are independent.
\end{enumerate}
\end{theorem}
\begin{proof}
\begin{board}[(a) Row \(i\) of \(K\) is \(t_i'X\), \(t_i\in\R^n\) (Theorem~\ref{h-thm-estimable}(a)); \(T\in\R^{q\times n}\)
has the rows \(t_i'\), so \(K=TX\) (\lec{1}, Theorem~4(c)); \(y\sim N_n(X\beta,\sigma^2I_n)\)
(Definition~\ref{h-def-normal}):]
\boardtop
\begin{align*}
K\beta^0&=TXGX'y=THy\sim N_q(THX\beta,\ \sigma^2THH'T')&&\text{(\lec{10}, Theorem~3)}\\
THX\beta&=TX\beta=K\beta,\qquad THH'T'=THT'=TXGX'T'=KGK'&&\text{(\lec{11}, Theorem~10(c))}
\end{align*}
\end{board}
\(H\) does not depend on \(G\) (\lec{11}, Theorem~10(b)), so neither do \(K\beta^0=THy\) and \(KGK'=THT'\). Under (A4),
\(X'X\) is invertible and \((X'X)^{-1}\) is its only g-inverse (\lec{11}, Theorem~2(b), Definition~8), every \(a'\beta\)
is estimable (Theorem~\ref{h-thm-estimable}(a)), and \(K=I_p\) gives \(\hat\beta\).
\begin{board}[(b) \(y\sim N_n(X\beta,\sigma^2I_n)\); \lec{11}, Theorem~10(c)--(e); \(z\sim N_n(\mu,\sigma^2I_n)\),
\(A'=A=A^2\), \(\rank A=d\geqslant1\), \(A\mu=0\) \(\Rightarrow\) \(z'Az/\sigma^2\sim\chi^2_d\) (\lec{10}, Theorem~15):]
\boardtop
\begin{align*}
\mathrm{SSE}&=y'My,\quad M'=M=M^2\\
\rank M&=n-r\geqslant1,\quad MX\beta=0\\
\mathrm{SSE}/\sigma^2&\sim\chi^2_{n-r}&&\text{(\lec{10}, Theorem~15, \(z=y\), \(A=M\))}\\
\prb(\mathrm{SSE}=0)&=\prb(\mathrm{SSE}/\sigma^2=0)=0&&\text{(\ref{f-quantile})}
\end{align*}
\end{board}
\begin{board}[(c) \(K\beta^0=(TH)y\) ((a)); for \(z\sim N_n(\mu,\sigma^2I_n)\), \(A'=A=A^2\) and \(KA=0\), \(Kz\) and
\(z'Az\) are independent (\lec{10}, Theorem~16(b)); here \(z=y\), and \(TH\), \(M\) for \(K\), \(A\):]
\[
(TH)M=T(HM)=0\implies K\beta^0\text{ and }\mathrm{SSE}=y'My\text{ independent}\qquad\text{(\lec{11}, Theorem~10(c))}
\]
\end{board}
\nopagebreak
\end{proof}
We write \(s^2=\hat\sigma^2=\mathrm{SSE}/(n-r)\), unbiased for \(\sigma^2\) (\lec{11}, Theorem~16:
\(\expc(\mathrm{SSE})=(n-r)\sigma^2\)), and \(s=\sqrt{s^2}\); by (b), \(s>0\) with probability \(1\).

\Needspace{18\baselineskip}
\begin{theorem}\label{h-thm-t}
In the normal linear model let \(\rank X=r<n\), \(G\) a g-inverse of \(X'X\), \(\beta^0=GX'y\), and let \(a\in\R^p\) be
nonzero with \(a'\beta\) estimable. Then \(a'Ga>0\); let \(\SE(a'\beta^0)=s\sqrt{a'Ga}\), the \textbf{standard error} of
\(a'\beta^0\), and \(t^*=t_{n-r,1-\alpha/2}\).
\begin{enumerate}
\item \(T=(a'\beta^0-a'\beta)/\SE(a'\beta^0)\sim t_{n-r}\).
\item \(\bigl[a'\beta^0-t^*\SE(a'\beta^0),\ a'\beta^0+t^*\SE(a'\beta^0)\bigr]\) is a confidence interval of level
\(1-\alpha\) for \(a'\beta\). For \(m_0\in\R\), the test of \(H_0\colon a'\beta=m_0\) that rejects when
\(\lvert T_0\rvert>t^*\), where \(T_0=(a'\beta^0-m_0)/\SE(a'\beta^0)\), has level \(\alpha\); it rejects exactly when
\(m_0\) lies outside this interval.
\end{enumerate}
\(\SE(a'\beta^0)\) and the interval are the same for every g-inverse \(G\) of \(X'X\) (Theorem~\ref{h-thm-sampling}(a),
\(K=a'\); \(s\) by \lec{11}, Theorem~10(e)). Under (A4), \(G=(X'X)^{-1}\), \(\beta^0=\hat\beta\) and \(r=p\).
\end{theorem}
\begin{proof}
\begin{board}[\(a'Ga>0\): \(a=X'c\) (Theorem~\ref{h-thm-estimable}(a)); \(H=XGX'\):]
\boardtop
\begin{align*}
a'Ga&=c'XGX'c=c'Hc=c'H'Hc=\lVert Hc\rVert^2&&\text{(\lec{11}, Theorem~10(c))}\\
X'Hc&=X'c=a\ne0\implies Hc\ne0\implies a'Ga>0&&\text{(\lec{6}, Definition~1)}
\end{align*}
\end{board}
\begin{board}[(a) \lec{10}, Theorem~3 for \(z\); \(s>0\) with probability \(1\):]
\boardtop
\begin{align*}
z&=\frac{a'\beta^0-a'\beta}{\sigma\sqrt{a'Ga}}\sim N(0,1)&&\text{(Theorem~\ref{h-thm-sampling}(a), \(K=a'\))}\\
v&=\mathrm{SSE}/\sigma^2\sim\chi^2_{n-r};\quad z,\ v\text{ independent}&&\text{(Theorem~\ref{h-thm-sampling}(b),(c); \ref{f-prob})}\\
T&=\frac z{s/\sigma}=\frac z{\sqrt{v/(n-r)}}\sim t_{n-r}&&\text{(\ref{f-chisq})}
\end{align*}
\end{board}
\begin{board}[(b) For \(T\sim t_\nu\): \(\prb(-t_{\nu,1-\alpha/2}\leqslant T\leqslant t_{\nu,1-\alpha/2})=1-\alpha\)
(\lec{10}, Theorem~11(a)) and \(t_{\nu,1-\alpha/2}>0\) (\lec{10}, Theorem~11(b)); \(\SE=\SE(a'\beta^0)>0\) with
probability \(1\):]
\boardtop
\begin{align*}
\prb(\lvert T\rvert\leqslant t^*)&=1-\alpha\quad(T\sim t_{n-r}\text{ by (a)})&&\text{(\lec{10}, Theorem~11(a))}\\
\lvert T\rvert\leqslant t^*&\iff a'\beta^0-t^*\SE\leqslant a'\beta\leqslant a'\beta^0+t^*\SE&&\text{(Definition~\ref{h-def-normal})}\\
H_0&\colon a'\beta=m_0\implies T_0=T\\
\prb(\lvert T_0\rvert>t^*)&=1-(1-\alpha)=\alpha\ \text{under }H_0&&\text{(\ref{f-prob})}\\
\lvert T_0\rvert>t^*&\iff\lvert a'\beta^0-m_0\rvert>t^*\SE\iff m_0\text{ outside the interval}
\end{align*}
\end{board}
\nopagebreak
\end{proof}

\Needspace{12\baselineskip}
\begin{example}\label{h-exm-t}
In Example~\ref{h-exm-estimable}, assume the normal linear model. Find the confidence interval of level \(1-\alpha\) for
\(\alpha_1-\alpha_2\), and the test of \(H_0\colon\alpha_1=\alpha_2\) at level \(\alpha\).
\end{example}
\textit{Solution.}
\begin{board}[\(\alpha_1-\alpha_2\) is estimable (Example~\ref{h-exm-estimable}); \(\beta^0\), \(a'Ga\) as in
Example~\ref{h-exm-solution}; \(\mathrm{SSE}=4\) (\lec{11}, Example~11), \(r=\rank X=3\) (\lec{11}, Theorem~6(b));
Theorem~\ref{h-thm-t}(b):]
\boardtop
\begin{align*}
a&=(0,1,-1,0)':\quad a'\beta^0=6-4=2,\quad a'Ga=\tfrac13+\tfrac12=\tfrac56;\quad\mathrm{SSE}=4,\ n-r=6-3=3,\ s^2=\tfrac43\\
\SE&=s\sqrt{5/6}=\sqrt{10/9}=\sqrt{10}/3;\quad\text{interval }\bigl[2-t_{3,1-\alpha/2}\sqrt{10}/3,\ 2+t_{3,1-\alpha/2}\sqrt{10}/3\bigr]\\
m_0&=0:\quad T_0=2/(\sqrt{10}/3)=6/\sqrt{10};\quad\text{the test rejects }H_0\iff t_{3,1-\alpha/2}<6/\sqrt{10}
\end{align*}
\end{board}
Equivalently, the interval contains \(0\) exactly when \(t_{3,1-\alpha/2}\geqslant6/\sqrt{10}\).

\subsection*{3. Testable hypotheses and the \(F\) test}

\begin{definition}\label{h-def-testable}
Let \(\rank X=r\), \(K\in\R^{q\times p}\) and \(m\in\R^q\). The hypothesis \(H_0\colon K\beta=m\) is \textbf{testable} if
every entry of \(K\beta\) is estimable and \(\rank K=q\) (the rows of \(K\) are linearly independent).
\end{definition}
\begin{board}[Let \(H_0\colon K\beta=m\) be testable, \(G\) a g-inverse of \(X'X\), \(H=XGX'\); \(K=T_1X\),
\(T_1\in\R^{q\times n}\) (Theorem~\ref{h-thm-estimable}(a); \lec{1}, Theorem~4(c)); for \(A\in\R^{n\times p}\), \(A'A\)
is PD (\lec{8}, Definition~13: symmetric, \(x'Ax>0\) for \(x\ne0\)) \(\iff\) \(\rank A=p\) (\lec{8}, Theorem~15(b)), and
a PD matrix is invertible (\lec{8}, Theorem~14(c)):]
\boardtop
\begin{align*}
T&=T_1H:\quad TX=T_1HX=T_1X=K,\quad TH=T_1H^2=T&&\text{(\lec{11}, Theorem~10(c))}\\
q&=\rank(TX)\leqslant\rank T\leqslant q\implies\rank T'=\rank T=q&&\text{(\lec{4}, Theorems~15(a), 12(a))}\\
V&=KGK'=TXGX'T'=THT'=THH'T'=TT'&&\text{(\lec{11}, Theorem~10(c))}\\
&\implies V\text{ is PD (\(A=T'\)), hence invertible}&&\text{(\lec{8}, Theorems~15(b), 14(c))}
\end{align*}
\end{board}
\(T=T_1H\) and \(V=THT'\) do not depend on \(G\) (\lec{11}, Theorem~10(b)). Under (A4), \(H_0\colon K\beta=m\) is testable
for every \(K\) of rank \(q\) (Theorem~\ref{h-thm-estimable}(a)).

\Needspace{22\baselineskip}
\begin{theorem}\label{h-thm-f}
In the normal linear model let \(\rank X=r<n\), \(G\) a g-inverse of \(X'X\), \(\beta^0=GX'y\), and \(H_0\colon K\beta=m\)
testable, with \(K\in\R^{q\times p}\) and \(V=KGK'\).
\begin{enumerate}
\item Let \(F=\bigl[(K\beta^0-m)'V^{-1}(K\beta^0-m)/q\bigr]/s^2\), defined with probability \(1\) (\(V\) invertible by
the lines after Definition~\ref{h-def-testable}, \(s^2>0\) by Theorem~\ref{h-thm-sampling}(b)) and the same for every
g-inverse \(G\) of \(X'X\). If \(H_0\) holds, then \(F\sim F_{q,n-r}\); so the test that rejects \(H_0\) when
\(F>F_{q,n-r,1-\alpha}\) has level \(\alpha\).
\item If \(q=1\) and \(K=a'\), then \(F=T_0^2\), with \(T_0\) of Theorem~\ref{h-thm-t}(b) for \(m_0=m\); the test of (a)
rejects exactly when that of Theorem~\ref{h-thm-t}(b) does.
\end{enumerate}
\end{theorem}
\begin{proof}
(a) \(K\beta^0\) and \(V\) do not depend on \(G\) (Theorem~\ref{h-thm-sampling}(a), the lines after
Definition~\ref{h-def-testable}), nor does \(s^2=y'My/(n-r)\) (\lec{11}, Theorem~10(b),(e)).
\begin{board}[Algebra: \(T\) as after Definition~\ref{h-def-testable} (\(TH=T\), \(V=TT'\)); \(H_0\) holds;
\(\varepsilon=y-X\beta\); \(MX=0\), \(X'M=0\) (\lec{11}, Theorem~10(c)); \(A'=A=A^2\Rightarrow\rank A=\tr A\) (\lec{8},
Theorem~11):]
\boardtop
\begin{align*}
K\beta^0&=TXGX'y=THy=Ty\\
K\beta^0-m&=Ty-TX\beta=T\varepsilon&&(m=K\beta=TX\beta)\\
U&=T'V^{-1}T=U',\quad U^2=T'V^{-1}(TT')V^{-1}T=U&&\text{(\lec{2}, Corollary~12)}\\
\rank U&=\tr U=\tr(V^{-1}TT')=\tr I_q=q&&\text{(\lec{1}, Theorem~10(b))}\\
TM&=T-TH=0\implies UM=0\\
Q&=(K\beta^0-m)'V^{-1}(K\beta^0-m)=\varepsilon'U\varepsilon;\ \mathrm{SSE}=y'My=\varepsilon'M\varepsilon
\end{align*}
\end{board}
\begin{board}[Distribution: \(\varepsilon\sim N_n(0,\sigma^2I_n)\) (A5); \(M'=M=M^2\), \(\rank M=n-r\) (\lec{11},
Theorem~10(c),(d)); for \(z\sim N_n(\mu,\sigma^2I_n)\) and symmetric \(A\), \(B\) with \(A^2=A\), \(B^2=B\), \(AB=0\),
\(z'Az\) and \(z'Bz\) are independent (\lec{10}, Theorem~16(a)):]
\boardtop
\begin{align*}
Q/\sigma^2&=\varepsilon'U\varepsilon/\sigma^2\sim\chi^2_q&&\text{(\lec{10}, Theorem~15, \(\mu=0\))}\\
\mathrm{SSE}/\sigma^2&=\varepsilon'M\varepsilon/\sigma^2\sim\chi^2_{n-r}&&\text{(\lec{10}, Theorem~15, \(\mu=0\))}\\
Q&\text{ and }\mathrm{SSE}\text{ independent}&&\text{(\lec{10}, Theorem~16(a), \(UM=0\))}\\
F&=\frac{(Q/\sigma^2)/q}{(\mathrm{SSE}/\sigma^2)/(n-r)}\sim F_{q,n-r}&&\text{(\ref{f-prob}; \ref{f-chisq})}\\
\prb(F&>F_{q,n-r,1-\alpha})=1-(1-\alpha)=\alpha&&\text{(\ref{f-quantile}, \ref{f-prob})}
\end{align*}
\end{board}
(b) \(K=a'\), \(V=a'Ga\), \(\SE(a'\beta^0)^2=s^2a'Ga\) (Theorem~\ref{h-thm-t}):
\[
F=\frac{(a'\beta^0-m)^2}{s^2a'Ga}=\Bigl(\frac{a'\beta^0-m}{\SE(a'\beta^0)}\Bigr)^2=T_0^2
\]
and \(F>F_{1,n-r,1-\alpha}\iff T_0^2>(t^*)^2\iff\lvert T_0\rvert>t^*\) (\lec{10}, Theorem~11(b):
\(t_{\nu,1-\alpha/2}^2=F_{1,\nu,1-\alpha}\), \(t_{\nu,1-\alpha/2}>0\)).
\end{proof}

\Needspace{19\baselineskip}
\begin{example}\label{b14.15}
Eight plots receive a fertilizer at two levels, coded \(x_{i1}=-1,1\), and irrigation at two levels, coded
\(x_{i2}=-1,1\), two plots for each combination; the yields are
\[
\begin{array}{c|rrrrrrrr}
x_{i1}&-1&-1&-1&-1&1&1&1&1\\
x_{i2}&-1&-1&1&1&-1&-1&1&1\\\hline
y_i&9&3&9&7&11&13&11&17
\end{array}
\]
In the normal linear model \(y_i=\beta_0+\beta_1x_{i1}+\beta_2x_{i2}+\varepsilon_i\) (multiple regression with \(n=8\),
\(k=2\), \lec{11}, Definition~1), find the confidence interval of level \(1-\alpha\) for \(\beta_1\), and the tests of
\(\beta_2=0\) and of \(\beta_1=\beta_2=0\) at level \(\alpha\). Put \(t^*=t_{5,1-\alpha/2}\).
\end{example}
\textit{Solution.}
\begin{board}[The three columns of \(X\) are orthogonal with squared length \(8\) (\lec{1}, Theorem~4(d));
\(X'X\) invertible \(\iff\) \(\rank X=3\) (\lec{11}, Theorem~2(b)); under (A4), \(G=(X'X)^{-1}\), \(\beta^0=\hat\beta\),
\(r=3\) (Theorem~\ref{h-thm-t}); every \(a'\beta\) is estimable (Theorem~\ref{h-thm-estimable}(a)):]
\boardtop
\begin{align*}
X'X&=8I_3,\quad\rank X=3,\quad G=(X'X)^{-1}=\tfrac18I_3,\quad X'y=(80,24,8)',\quad\hat\beta=\tfrac18X'y=(10,3,1)'\\
\hat y_i&=10+3x_{i1}+x_{i2}=6,\ 8,\ 12,\ 14\ \text{for the four combinations}\\
e&=y-\hat y=(3,-3,1,-1,-1,1,-3,3)',\quad\mathrm{SSE}=e'e=40,\quad s^2=40/(8-3)=8
\end{align*}
\end{board}
\begin{board}[\(\SE(\hat\beta_j)=s\sqrt{1/8}=1\) (\(a'Ga=\frac18\) for \(a=\col_{j+1}(I_3)\)); Theorems~\ref{h-thm-t}(b),
\ref{h-thm-f}(b):]
\boardtop
\begin{align*}
\beta_1&:\ \text{interval }[3-t^*,\ 3+t^*]\\
\beta_2&=0:\ T_0=(1-0)/1=1;\ \text{the test rejects}\iff t^*<1;\quad F=T_0^2=1
\end{align*}
\end{board}
\begin{board}[\(\beta_1=\beta_2=0\) is \(K\beta=0\), \(K=\left(\begin{smallmatrix}0&1&0\\0&0&1\end{smallmatrix}\right)\),
\(\rank K=2\): testable under (A4) (the lines after Definition~\ref{h-def-testable}); Theorem~\ref{h-thm-f}(a),
\(q=2\):]
\boardtop
\begin{align*}
K\hat\beta&=(3,1)',\quad V=KGK'=\tfrac18I_2,\quad V^{-1}=8I_2,\quad(K\hat\beta)'V^{-1}K\hat\beta=8(9+1)=80\\
F&=\frac{80/2}8=5:\ \text{the test rejects }H_0\iff F_{2,5,1-\alpha}<5
\end{align*}
\end{board}
With \(\bar y=10\): \(\mathrm{SSR}=2(16+4+4+16)=80\), \(\mathrm{SST}=120\) and \(R^2=\frac23\) (\lec{11}, Definition~17:
\(\mathrm{SSR}=\sum_i(\hat y_i-\bar y)^2\), \(\mathrm{SST}=\sum_i(y_i-\bar y)^2\), \(R^2=\mathrm{SSR}/\mathrm{SST}\)); so
here \((K\hat\beta)'V^{-1}K\hat\beta=80=\mathrm{SSR}\).

\Needspace{20\baselineskip}
\subsection*{4. One-way analysis of variance}

\Needspace{20\baselineskip}
\begin{theorem}\label{h-thm-oneway-anova}
In the one-way design (\lec{11}, Theorem~6: \(n\) observations in \(g\) groups of sizes \(n_l\geqslant1\), the
indicator \(d_l\in\R^n\) of group \(l\), \(X=[\one_n\ d_1\ \cdots\ d_g]\), \(X_1=[d_1\ \cdots\ d_g]\)), let
\(\mathcal I_l\) be the set of the observations (row indices) of group \(l\), \(\bar y_l=n_l^{-1}\sum_{i\in\mathcal
I_l}y_i\) the \textbf{group mean} and \(\bar y\) the \textbf{grand mean} of all \(n\) observations. Then
\(\hat y=\bar y_1d_1+\dots+\bar y_gd_g\) (the fitted value of each observation is the mean of its group), and
\[
\sum_{i=1}^n(y_i-\bar y)^2=\sum_{l=1}^gn_l(\bar y_l-\bar y)^2+\sum_{l=1}^g\sum_{i\in\mathcal I_l}(y_i-\bar y_l)^2
\]
that is, \(\mathrm{SST}=\mathrm{SSR}+\mathrm{SSE}\), with the \textbf{between-group} sum of squares \(\mathrm{SSR}\) and the
\textbf{within-group} sum of squares \(\mathrm{SSE}\).
\end{theorem}
\begin{proof}
\begin{board}[Fitted values: \(X_1\) has full column rank (\lec{11}, Theorem~6(a)); \(X_1'y\) holds the group totals
\(d_l'y\) (\lec{1}, Theorem~4(d)); no observation lies in two groups; \(\diag(n_l)^{-1}=\diag(1/n_l)\) (\lec{1},
Proposition~13(b); \lec{2}, Theorem~15):]
\boardtop
\begin{align*}
d_l'd_m&=0\ (l\ne m),\quad d_l'd_l=n_l\implies X_1'X_1=\diag(n_1,\dots,n_g)&&\text{(\lec{1}, Theorem~4(d))}\\
\hat\beta_1&=(X_1'X_1)^{-1}X_1'y=(\bar y_1,\dots,\bar y_g)'&&\text{(\lec{11}, Theorem~3(c))}\\
X_1\hat\beta_1&=\bar y_1d_1+\dots+\bar y_gd_g=\hat y\quad(\mathcal C(X_1)=\mathcal C(X))&&\text{(\lec{11}, Theorems~6(c), 3(b))}
\end{align*}
\end{board}
\begin{board}[Sums of squares: \(\hat y_i=\bar y_l\) for \(i\in\mathcal I_l\):]
\boardtop
\begin{align*}
\mathrm{SSR}&=\textstyle\sum_i(\hat y_i-\bar y)^2=\sum_ln_l(\bar y_l-\bar y)^2\\
\mathrm{SSE}&=\textstyle\sum_i(y_i-\hat y_i)^2=\sum_l\sum_{i\in\mathcal I_l}(y_i-\bar y_l)^2&&\text{(\lec{11}, Definition~17)}\\
\mathrm{SST}&=\mathrm{SSR}+\mathrm{SSE}\quad(\one_n=\textstyle\sum_ld_l\in\mathcal C(X))&&\text{(\lec{11}, Theorem~18)}
\end{align*}
\end{board}
\nopagebreak
\end{proof}

\Needspace{22\baselineskip}
\begin{corollary}\label{h-cor-oneway-f}
In the one-way design of \lec{11}, Theorem~6, use the design \(X_1=[d_1\ \cdots\ d_g]\) with \(\beta=(\mu_1,\dots,\mu_g)'\)
(the group means), and assume the normal linear model with \(n>g\geqslant2\). With \(\mathrm{SSR}\), \(\mathrm{SSE}\) of
Theorem~\ref{h-thm-oneway-anova}, let \(F=\bigl[\mathrm{SSR}/(g-1)\bigr]/\bigl[\mathrm{SSE}/(n-g)\bigr]\), defined with
probability \(1\) (Theorem~\ref{h-thm-sampling}(b)). If \(H_0\colon\mu_1=\dots=\mu_g\) holds, then
\(F\sim F_{g-1,n-g}\); so the test that rejects \(H_0\) when \(F>F_{g-1,n-g,1-\alpha}\) has level \(\alpha\).
\end{corollary}
\begin{proof}
\(\bar J_n=\frac1n\one\one'\) (\lec{1}, Definition~14), \(\bar J_n'=\bar J_n=\bar J_n^2\), \(\tr\bar J_n=1\) (\lec{1},
Proposition~15), \(\bar J_ny=\bar y\one_n\) (\(\bar y=\frac1n\one'y\), \lec{1}, Example~16); \(\one_n=X_1\one_g\).
\begin{board}[\(\rank X_1=g\) (\lec{11}, Theorem~6(a)), so \(H=X_1(X_1'X_1)^{-1}X_1'\) and \(M=I_n-H\) have \(H'=H=H^2\),
\(HX_1=X_1\), \(HM=0\), \(\tr H=g\), \(Hy=\hat y\) (\lec{11}, Theorems~2(b), 10, design \(X_1\)); \(A=H-\bar J_n\):]
\boardtop
\begin{align*}
H\bar J_n&=\tfrac1n(HX_1\one_g)\one'=\tfrac1n\one\one'=\bar J_n&&\text{(\lec{1}, Theorem~3(a))}\\
\bar J_nH&=(H\bar J_n)'=\bar J_n&&\text{(\lec{1}, Theorem~8)}\\
A'&=A,\quad A^2=H-2\bar J_n+\bar J_n^2=A\\
\rank A&=\tr A=\tr H-\tr\bar J_n=g-1&&\text{(\lec{8}, Theorem~11; \lec{1}, Theorem~10(a))}\\
AM&=HM-\bar J_n(I_n-H)=-\bar J_n+\bar J_n=0\\
y'Ay&=y'A'Ay=\lVert\hat y-\bar y\one_n\rVert^2=\mathrm{SSR}&&\text{(Theorem~\ref{h-thm-oneway-anova})}\\
y'My&=e'e=\mathrm{SSE}&&\text{(\lec{11}, Theorem~10(e))}
\end{align*}
\end{board}
\begin{board}[Distribution: \(y\sim N_n(X_1\beta,\sigma^2I_n)\) (Definition~\ref{h-def-normal}); \(g-1\geqslant1\);
\(\bar J_n\one_n=\frac1n\one(\one'\one)=\one_n\) (\lec{1}, Proposition~15):]
\boardtop
\begin{align*}
AX_1\beta&=\mu_1A(d_1+\dots+d_g)=\mu_1(H\one_n-\bar J_n\one_n)=\mu_1(\one_n-\one_n)=0&&\text{(\(H_0\))}\\
\mathrm{SSR}/\sigma^2&=y'Ay/\sigma^2\sim\chi^2_{g-1}&&\text{(\lec{10}, Theorem~15)}\\
\mathrm{SSE}/\sigma^2&=y'My/\sigma^2\sim\chi^2_{n-g}&&\text{(Theorem~\ref{h-thm-sampling}(b), \(r=g\))}\\
\mathrm{SSR}&\text{ and }\mathrm{SSE}\text{ are independent }(AM=0)&&\text{(\lec{10}, Theorem~16(a))}\\
F&=\frac{(\mathrm{SSR}/\sigma^2)/(g-1)}{(\mathrm{SSE}/\sigma^2)/(n-g)}\sim F_{g-1,n-g},\ \prb(F>F_{g-1,n-g,1-\alpha})=\alpha&&\text{(\ref{f-prob}--\ref{f-quantile})}
\end{align*}
\end{board}
\nopagebreak
\end{proof}

\Needspace{10\baselineskip}
The computations of Corollary~\ref{h-cor-oneway-f} are arranged in the \textbf{analysis-of-variance table}: the degrees of
freedom are those of the \(\chi^2\) distributions in the proof, and \(n-1=(g-1)+(n-g)\) for \(\mathrm{SST}\); the
\textbf{mean squares} are the sums of squares divided by them, and \(F\) is their ratio:
\[
\begin{array}{lccc}
\text{source}&\text{degrees of freedom}&\text{sum of squares}&\text{mean square}\\\hline
\text{between groups}&g-1&\mathrm{SSR}=\sum_ln_l(\bar y_l-\bar y)^2&\mathrm{SSR}/(g-1)\\
\text{within groups}&n-g&\mathrm{SSE}=\sum_l\sum_{i\in\mathcal I_l}(y_i-\bar y_l)^2&s^2=\mathrm{SSE}/(n-g)\\\hline
\text{total}&n-1&\mathrm{SST}=\sum_i(y_i-\bar y)^2&
\end{array}
\]

\Needspace{10\baselineskip}
\begin{example}\label{h-exm-oneway-f}
In Example~\ref{h-exm-estimable}, assume the normal linear model. Test \(H_0\colon\alpha_1=\alpha_2=\alpha_3\) at level
\(\alpha\).
\end{example}
\textit{Solution.}
\begin{board}[Group means \(6,4,2\) (\lec{11}, Example~7); \(e=(-1,1,0,-1,1,0)'\), \(\mathrm{SSE}=4\) (\lec{11},
Example~11); Theorem~\ref{h-thm-oneway-anova}:]
\boardtop
\begin{align*}
\bar y&=28/6=14/3;\quad\mathrm{SSE}=4,\ r=\rank X=3,\ s^2=4/3\\
\mathrm{SSR}&=3\Bigl(6-\frac{14}3\Bigr)^2+2\Bigl(4-\frac{14}3\Bigr)^2+\Bigl(2-\frac{14}3\Bigr)^2=\frac{48+8+64}9=\frac{40}3,
\qquad\mathrm{SST}=\frac{40}3+4=\frac{52}3
\end{align*}
\end{board}
\begin{board}[\(\one_6=d_1+d_2+d_3\); Corollary~\ref{h-cor-oneway-f} with \(n=6>g=3\):]
\boardtop
\begin{align*}
X\beta&=\mu\one_6+\alpha_1d_1+\alpha_2d_2+\alpha_3d_3=X_1(\mu+\alpha_1,\ \mu+\alpha_2,\ \mu+\alpha_3)':\ H_0\text{ is }\mu_1=\mu_2=\mu_3\\
F&=\frac{(40/3)/2}{4/3}=5,\ \text{with }2\text{ and }3\text{ degrees of freedom: the test rejects }H_0\iff F_{2,3,1-\alpha}<5
\end{align*}
\end{board}
\[
\begin{array}{lccc}
\text{source}&\text{degrees of freedom}&\text{sum of squares}&\text{mean square}\\\hline
\text{between groups}&2&40/3&20/3\\
\text{within groups}&3&4&4/3\\\hline
\text{total}&5&52/3&
\end{array}
\]
\begin{board}[The same \(F\) by Theorem~\ref{h-thm-f}(a) in the design \(X\) (\(r=3\)): \(H_0\) is \(K\beta=0\), testable
(the criterion of Example~\ref{h-exm-estimable} for the rows); \(G\), \(\beta^0=(0,6,4,2)'\) as in
Example~\ref{h-exm-solution}; the rows of \(K\) are independent \(\iff\) \(\rank K=2\) (\lec{4}, Theorem~12(c));
\(V^{-1}\) by \lec{2}, Proposition~17:]
\boardtop
\begin{align*}
&K=\begin{pmatrix}0&1&-1&0\\0&0&1&-1\end{pmatrix}:\ -0+1-1+0=0,\ -0+0+1-1=0:\ \text{rows estimable}\\
&c_1(0,1,-1,0)+c_2(0,0,1,-1)=(0,\ c_1,\ c_2-c_1,\ -c_2)=0\implies c_1=c_2=0:\ \rank K=2\\
&K\beta^0=(2,2)',\quad V=KGK'=\begin{pmatrix}\frac56&-\frac12\\-\frac12&\frac32\end{pmatrix}\\
&ad-bc=\tfrac54-\tfrac14=1,\quad V^{-1}=\begin{pmatrix}\frac32&\frac12\\\frac12&\frac56\end{pmatrix}\\
&(K\beta^0)'V^{-1}K\beta^0=4\bigl(\tfrac32+\tfrac12+\tfrac12+\tfrac56\bigr)=\tfrac{40}3=\mathrm{SSR},\qquad
F=\frac{40/3}{2\cdot4/3}=5
\end{align*}
\end{board}
Uncentered, \(y'y=n\bar y^2+\mathrm{SSR}+\mathrm{SSE}\) (\lec{11}, Theorem~18): \(n\bar y^2=6\cdot\frac{196}9=\frac{392}3\),
and \(y'y=148=\frac{392}3+\frac{40}3+4\).

\bigskip
\textbf{Practice.} The Exercises and Problems of the textbook, Chapter~14: estimable functions and their best linear
unbiased estimators, the normal linear model, \(t\) intervals and tests, testable hypotheses and the \(F\) test, and the
one-way analysis of variance.

\end{document}
